\section{Experiments}
\label{sec:experiments}

\paragraph{Setup.}
We independently pre-train Panda V2 on PILArNet-M~\citep{polarmae}, TPCpp-10M~\citep{tpcpp10m}, and a preliminary version of the Water-Cherenkov Annotated Neutrino Dataset (WAND), representing LArTPC, collider-TPC, and SK-like water Cherenkov data, respectively. All downstream modules are trained using exactly \textbf{1,000 labeled events}. On PILArNet-M and TPCpp-10M, we perform semantic segmentation with a linear classifier using 414K trainable LoRA~\citep{lora} parameters, while particle clustering uses a 580K parameter two-stage GNN together with LoRA. On WAND, frozen lightweight readouts reconstruct single-particle properties, while Panda Detector with LoRA performs multi-ring reconstruction. For Panda V1~\citep{panda}, we compare against its strongest published frozen-backbone result for each task: Panda Detector for clustering and a 16M-parameter U-Net-like decoder for semantic segmentation.

\subsection{Reconstruction from 1,000 labeled events}
\label{sec:reconstruction_results}

Figure~\ref{fig:reconstruction} summarizes reconstruction performance across the three modalities. On PILArNet-M, Panda V2's GNN+LoRA reaches a particle-clustering ARI of $0.973$, compared with $0.914$ for Panda V1's frozen-backbone Panda Detector at 1,000 labels, and matches its $0.973$ result at one million labels. For semantic segmentation, linear+LoRA reaches a macro-F1 of $0.976$, compared with $0.952$ for Panda V1's 16M-parameter frozen-backbone decoder at 1,000 labels and $0.978$ at 100,000 labels.

On TPCpp-10M, the graph probe with LoRA reaches a track clustering ARI of $0.943$, closely matching FM4NPP's $0.945$ result obtained with 70,000 labeled events. A linear semantic head with LoRA reaches macro-F1 scores of $0.838$ for particle identification and $0.942$ for noise tagging, compared with $0.713$ and $0.912$ for FM4NPP, respectively.

For SK-like WAND data, simple probes from Panda V2 features reconstruct single-particle kinematics using only 1,000 labeled events, reaching electron (muon) vertex resolutions of $98.7$ (N/A) cm, angular resolutions of $5.75^\circ$ ($5.88^\circ$), and momentum resolutions of $4.80\%$ ($4.82\%$). These remain worse than the highly optimized fiTQun and APFit reconstruction algorithms, but demonstrate that useful detector-level observables are accessible through inexpensive adaptation of the same representation learner. For multi-ring events, Panda Detector obtains recalls of $80.4\%$, $44.4\%$, and $96.7\%$ for one-, two-, and $\geq3$-ring events, respectively. Together, these results show that a common pre-trained representation can support full reconstruction from only 1,000 labeled events.


\subsection{Physical structure in frozen representations}
\label{sec:interpretability}

\begin{figure}
    \centering
    \includegraphics[width=0.9\linewidth]{figs/interp.jpg}
    \caption{Single directions in the frozen Panda V2 representation recover physically meaningful structure. (a) In LArTPC events, a latent direction orders points along particle trajectories despite the absence of timing information. (b) In TPCpp-10M, a latent direction encodes track curvature and transverse momentum.}
    \label{fig:interp}
\end{figure}

We next ask whether the representations encode further physically meaningful quantities. We train a single vectors to predict physically meaningful quantities when projected onto per-point features. As shown in Fig.~\ref{fig:interp}, a single learned vector, when projected onto per point features with the inner product, produces a scalar field that orders points along particle trajectories in LArTPC events (Spearman $\rho=0.89$), recovering an emergent ``arrow of time'' despite the absence of timing information in the input. In TPCpp-10M, another latent direction encodes track curvature and thus transverse momentum (Spearman $\rho=0.957$). More information and examples can be found in the appendix.
